\section{The Goren--Oort operator}\label{sec:operator}

This section defines the matrix $A$ and relates its kernels to Tanner codes on covers of a finite graph. The arithmetic
that makes these kernels large is in \cref{sec:hypothesis}; here we need only the combinatorial structure.

\subsection{The icosian group and the base graph}\label{sec:icosian}

Let $E=\Q(\sqrt5)$, with ring of integers $\OO_E=\Z[\theta]$, $\theta=(1+\sqrt5)/2$. The prime $3$ is inert in~$E$; we
identify $\OO_E/3$ with $\F_9=\F_3[i]/(i^2+1)$ by $\theta\mapsto2+i$. The prime $11$ splits; let $\lambda_0$ be the prime
above~$11$ with $\theta\equiv4\pmod{\lambda_0}$, so that $\OO_E/\lambda_0^m\cong\Z/11^m$. Let $B$ be the quaternion algebra
over $E$ ramified exactly at the two infinite places, and $I\subset B$ the icosian ring, a maximal order of class
number~$1$ whose norm-one units form the binary icosahedral group $2I\cong\SL_2(\F_5)$ of order
$120$~\cite[Proposition~V.3.4 and p.~155]{Vigneras1980}. Since $B$ splits at~$3$, $B\otimes_EE_3\cong M_2(\Q_9)$, where $\Q_9$ is the unramified
quadratic extension of~$\Q_3$ with ring of integers $\Z_9$, and
\[
  \Gamma^+=I[1/3]^1/\{\pm1\}\subset\PSL_2(\Q_9)
\]
acts on the Bruhat--Tits tree $\Tree$ of $\PGL_2(\Q_9)$, which is $10$-regular; it preserves the bipartition of the
vertices into even and odd ones. Let $\Lambda_0$ be the vertex fixed by $I\otimes\OO_{E,3}\cong M_2(\Z_9)$, let
$\alpha=1+i+j\in I$, of reduced norm~$3$, and put $\Lambda_1=\alpha\Lambda_0$, a neighbour of~$\Lambda_0$. The stabilizers
of $\Lambda_0$, $\Lambda_1$ and the edge $e_0=\{\Lambda_0,\Lambda_1\}$ are $A_5=2I/\{\pm1\}$, $\alpha A_5\alpha^{-1}$ and a
group $S_3$, whose preimage $\Gamma_e\subset2I$ is binary dihedral of order~$12$; and
$\Gamma^+=A_5*_{S_3}\alpha A_5\alpha^{-1}$ (\cref{lem:amalgam}).

For $m\ge1$ let $\Gamma^+(\lambda_0^m)=\ker\bigl(\Gamma^+\to\PSL_2(\OO_E/\lambda_0^m)\bigr)$. It is torsion-free and acts
freely on~$\Tree$ (\cref{lem:chain}). The \emph{base graph} is $X_0=\Gamma^+(\lambda_0)\backslash\Tree$.

\begin{lemma}[The base graph]\label{lem:basegraph}
\notformalized
Let $P=\PSL_2(\F_{11})$, and let $H_0$, $H_1$ and $H_e$ be the images of $A_5$, $\alpha A_5\alpha^{-1}$ and $S_3$ under
reduction modulo~$\lambda_0$. Then $H_e=H_0\cap H_1$, and $X_0$ is the coset graph with even vertices $P/H_0$, odd
vertices $P/H_1$ and edges $P/H_e$, the edge $gH_e$ joining $gH_0$ to $gH_1$. It has $11$ vertices of each parity and
$110$ edges, every even vertex is joined to $10$ distinct odd vertices, and $X_0$ is the complete bipartite graph
$K_{11,11}$ with a perfect matching removed. Its fundamental group is free of rank $110-22+1=89$.
\end{lemma}
\begin{informalproof}
By \cref{lem:chain}, $\Gamma^+\to P$ is onto with kernel $\Gamma^+(\lambda_0)$, which meets no vertex stabilizer. Since
$\Gamma^+$ acts transitively on the even vertices, on the odd vertices and on the edges (\cref{lem:amalgam}), the even
vertices of $X_0$ are $\Gamma^+(\lambda_0)\backslash\Gamma^+/A_5=P/H_0$, and likewise for odd vertices and edges, with
incidences as stated. This gives $660/60=11$ vertices of each parity and $660/6=110$ edges. The edges $hH_e$
($h\in H_0$) at the even vertex $H_0$ have odd ends $hH_1$, and two of them have the same odd end exactly when they differ
by an element of $H_0\cap H_1$. This group contains $H_e\cong S_3$, which is maximal in $H_0\cong A_5$. If it were all of
$H_0$, then $H_1=H_0$, every edge at $gH_0$ would end at the odd vertex $gH_1$, and the quotient $X_0$ of the connected
tree $\Tree$ would be disconnected. So $H_0\cap H_1=H_e$, and the ten edges at an even vertex end at ten distinct odd
vertices. By symmetry every odd vertex is joined to ten distinct even vertices, so the missing pairs form a perfect
matching. A connected graph with $22$ vertices and $110$ edges has free fundamental group of rank $89$.
\end{informalproof}

Lifting a spanning tree of $X_0$ to $\Tree$ identifies $\pi_1(X_0)$ with $\Gamma^+(\lambda_0)$; we write $\psi$ for this
isomorphism (\cref{lem:covering}).

\subsection{Local codes and data}\label{sec:datum}

The $10$ edges at a vertex $\Lambda$ correspond to the lattices $\Lambda'$ with $3\Lambda\subset\Lambda'\subset\Lambda$
of index~$9$, that is, to the points of the projective line $\Pone(\Lambda/3\Lambda)$ over~$\F_9$.

\begin{definition}[Local codes]\label{def:localcode}
For vectors $x_1,\dots,x_{10}\in\F_9^2$ representing the $10$ points of $\Pone(\F_9)$, the \emph{local code} is
$C=\{(f(x_1),\dots,f(x_{10})): f\text{ a binary quartic form over }\F_9\}\subseteq\F_9^{10}$.
\end{definition}

\begin{lemma}[Self-duality]\label{lem:selfdual}
\notformalized
The local code has dimension $5$ and equals its dual for the standard bilinear form, for every choice of the
representatives $x_j$.
\end{lemma}
\begin{informalproof}
A nonzero binary quartic vanishes at no more than four points of $\Pone$, so evaluation is injective and $\dim C=5$. For
quartics $f,g$ put $h=fg$, of degree~$8$. Since $c^8=1$ for $c\in\F_9^\times$, $h(cx)=h(x)$, so
$\sum_{v\in\F_9^2\setminus0}h(v)=8\sum_jh(x_j)=-\sum_jh(x_j)$. For a monomial $h=X^aY^{8-a}$,
$\sum_{(s,t)\in\F_9^2}s^at^{8-a}=\bigl(\sum_ss^a\bigr)\bigl(\sum_tt^{8-a}\bigr)$, and $\sum_{s\in\F_9}s^a=0$ unless $a$ is a
positive multiple of~$8$ (with $s^0=1$, the sum for $a=0$ is $9=0$). One of the two factors therefore vanishes for every
$0\le a\le8$, and $\sum_jf(x_j)g(x_j)=0$. So $C\subseteq C^\perp$, and equality holds because $\dim C=5=10-\dim C$.
\end{informalproof}

A \emph{datum} on $\Tree$, or on a quotient of $\Tree$, assigns to every vertex a local code on the edges at that vertex,
given by labels in the sense of \cref{def:localcode}. Its \emph{Tanner code} on a quotient graph $X$ consists of the
functions on the edges of $X$ whose restriction to the edges at each vertex lies in the local code there.

Choose $p_0\in\F_9^2$ spanning the image and $k_0$ spanning the kernel of $\alpha\bmod3\in M_2(\F_9)$, in the coordinates
of $\Lambda_0/3\Lambda_0=\F_9^2$. The edges at $\Lambda_0$ are $u\,e_0$ for $u\in2I$; label $u\,e_0$ by the vector
$u\,p_0$. At $\Lambda_1$, in the coordinates transported by~$\alpha$, label the edge $t\,e_0$, where
$t=\alpha t'\alpha^{-1}$ with $t'\in2I$, by $\Frob(t'k_0)$, where $\Frob$ cubes both coordinates. Two elements give the
same edge exactly when they differ by an element of~$\Gamma_e$, and then their labels differ by a scalar $c$ with
$c^4=1$: $\Gamma_e$ acts on the line of the edge through a character, which factors through the abelianization $C_4$ of
the binary dihedral group of order~$12$. The space of binary quartic forms is $\GL_2(\F_9)$-stable, and $f(cx)=f(x)$ when
$c^4=1$; so the local codes at $\Lambda_0$ and $\Lambda_1$ are well defined, and transporting them by $\Gamma^+$ gives a
$\Gamma^+$-equivariant datum on $\Tree$, the \emph{twisted datum} $D_{\mathrm{tw}}$. The \emph{untwisted datum}
$D_{\mathrm{untw}}$ is the same without $\Frob$. For $\Gamma\subseteq\Gamma^+(\lambda_0)$ these data descend to
$\Gamma\backslash\Tree$; write $T_{\mathrm{tw}}(\Gamma\backslash\Tree)$ for the Tanner code of the descent of
$D_{\mathrm{tw}}$. \Cref{sec:hypothesis} shows that the datum that geometry produces has the type of $D_{\mathrm{tw}}$ or
of its image under the Frobenius (\cref{lem:twist}), which have Tanner codes of the same dimensions.

\subsection{The matrix}\label{sec:matrix}

Fix a spanning tree $\tau$ of~$X_0$. Its $21$ edges leave $89$ edges, and $F=\pi_1(X_0,\tau)$ is free on generators
$x_0,\dots,x_{88}$, one for each edge outside~$\tau$, in the order of the data. Put $\gamma_e$ equal to the generator of
$e$ for an edge outside $\tau$ and $\gamma_e=1$ for an edge of~$\tau$. For an edge~$e$, let
$m_0(e)=(x^{4-a}y^a)_{a=0,\dots,4}$ be the quartic monomials evaluated at the label of $e$ at its even end, and $m_1(e)$
the same at its odd end, with the twisted labels.

\begin{definition}[The operator]\label{def:operator}
\formal{LuckApproximation/Operator.lean}{Luck.Op.F,Luck.Op.gamma,Luck.Op.entryEdge,Luck.Op.coupling,Luck.Op.opA}
The $55\times55$ matrix $A_9$ over $\F_9[F]$ has rows indexed by pairs $(v_0,c)$ of an even vertex and $c\in\{0,\dots,4\}$
and columns by pairs $(v_1,r)$ of an odd vertex and $r\in\{0,\dots,4\}$, and
\[
  A_9[(v_0,c),(v_1,r)]=\sum_{e\colon v_0\text{---}v_1}m_1(e)_r\,m_0(e)_c\,\gamma_e\in\F_9[F].
\]
By \cref{lem:basegraph} each entry has at most one term. The matrix $A\in M_{110}(\F_3[F])$ is the restriction of scalars
of $A_9$ along $\F_9=\F_3[i]/(i^2+1)$: an entry $c_0+c_1i$ becomes the block
$\bigl(\begin{smallmatrix}c_0&c_1\\-c_1&c_0\end{smallmatrix}\bigr)$, row $(v_0,c)$, component~$a$, has index
$10v_0+2c+a$, and column $(v_1,r)$, component~$b$, has index $10v_1+2r+b$. For a field $K$ of characteristic~$3$, $A_K$ is the image of $A$ under $\F_3\to K$ (\lean{Luck.Op.opK}).
\end{definition}

In Lean the matrix is this construction, computed on natural numbers (\path{LuckApproximation/Construction.lean}).
The $120$ icosians are given by the standard formula (\lean{Luck.Con.icosians}): the units $\pm1,\pm i,\pm j,\pm k$, the
elements $(\pm1\pm i\pm j\pm k)/2$, and the elements whose coordinates are an even permutation of
$(0,\pm\frac12,\pm\frac{\theta^{-1}}2,\pm\frac\theta2)$. They are reduced by the splittings $i\mapsto(\begin{smallmatrix}1&1\\1&-1\end{smallmatrix})$,
$j\mapsto(\begin{smallmatrix}0&1\\-1&0\end{smallmatrix})$ over $\F_9$ and $i\mapsto(\begin{smallmatrix}1&3\\3&-1\end{smallmatrix})$,
$j\mapsto(\begin{smallmatrix}0&1\\-1&0\end{smallmatrix})$ over $\F_{11}$ (\lean{Luck.Con.red3}, \lean{Luck.Con.red11}), and the
labels, the quartic monomials and the blocks are computed from them (\lean{Luck.Con.label0}, \lean{Luck.Con.label1},
\lean{Luck.Con.mono}, \lean{Luck.Con.blockNat}). Generated data enter only as a coset certificate and packed caches,
both checked by the kernel, and as free choices. The certificate (\path{ConData.lean}) gives coset representatives of the
vertices, the position icosians, and the positions of each edge at its two ends; \cref{lem:enumeration} checks it. Packed
copies of the incidences and blocks speed up the certificate of \cref{sec:certificate}; the kernel checks that they agree
with the construction (\lean{Luck.Con.blockNat_eq_cache}, \lean{Luck.Con.edgeOfNat_eq_cache}), so nothing in the proof
trusts them. The choices are the spanning tree, with a numbering of the other edges, which \cref{lem:enumeration} checks
to be a spanning tree, and the voltage $\chi$ of \cref{sec:certificate}. The coset representatives are also a choice: they
fix coordinates at each vertex. Another choice multiplies the labels at a vertex by a fixed matrix in the image of $2I$,
and so multiplies its five rows (or columns) of $A_9$ by a fixed invertible $5\times5$ matrix over~$\F_9$, which changes no
kernel dimension; another choice of position icosians multiplies labels by scalars $c$ with $c^4=1$ and leaves the matrix
unchanged.

\begin{lemma}[The Lean enumeration]\label{lem:enumeration}
\formal{LuckApproximation/ConChecks.lean}{Luck.Con.splittings_ok,Luck.Con.subgroups_ok,Luck.Con.cosetGraph_ok,Luck.Con.spanningTree_ok,Luck.Con.labels_ok}
Let $\pm H_0$ and $\pm H_1$ be the sets of images in $\SL_2(\F_{11})$ of the $120$ icosians and of their conjugates
by~$\alpha$, and of their negatives, and $\pm H_e=\pm H_0\cap\pm H_1$. Then:
\begin{enumerate}[label=(\alph*)]
\item the two splittings respect $i^2=j^2=-1$ and $ij=-ji$, and $\alpha$ is invertible modulo~$\lambda_0$;
\item $\pm H_0$ and $\pm H_1$ consist of matrices of determinant~$1$, and the product of any two of the $120$ images lies
in $\pm H_0$ (resp.\ $\pm H_1$); so $\pm H_0$ and $\pm H_1$ are subgroups of $\SL_2(\F_{11})$;
\item the representatives of the even vertices have determinant~$1$ and lie in distinct cosets of $\pm H_0$ that cover
$\SL_2(\F_{11})$, and likewise the representatives of the odd vertices and $\pm H_1$;
\item the ten position icosians at the even vertices, and the conjugates by~$\alpha$ of the ten position icosians at the
odd vertices, represent the ten cosets of $\pm H_e$ in $\pm H_0$ (resp.\ $\pm H_1$) once each;
\item the ends and positions of the edges are in range, no two edges have the same ends, the pairs of an end and a
position are distinct, and each edge is the same coset of $\pm H_e$ seen from its two ends;
\item the edges without a generator form a spanning tree, and the other edges carry distinct generators;
\item $p_0\ne0$, $k_0\ne0$ and $(\alpha\bmod3)\,k_0=0$; since $p_0$ is a column of $\alpha\bmod3$, it spans the image
and $k_0$ the kernel.
\end{enumerate}
\end{lemma}
\begin{outline}
Each statement is a Boolean computation on natural numbers, evaluated by the kernel (\texttt{decide +kernel}). Matrices
over $\F_{11}$ are coded as single numbers and the sets $\pm H_0$, $\pm H_1$ as bit masks, which are closed under negation
by construction; coverage is checked over all $11^4$ matrices.
\end{outline}

By~(b), $\pm H_0$ and $\pm H_1$ are subgroups, the preimages in $\SL_2(\F_{11})$ of $H_0$ and $H_1$, and by
\cref{lem:basegraph} $\pm H_e$ is the preimage of~$H_e$. So (c)--(e) say that the enumeration is the coset graph of
\cref{lem:basegraph}, and~(g) that the labels are those of \cref{sec:datum}. So the Lean operator \lean{Luck.Op.opA} is
the matrix of \cref{def:operator}, with no step resting on generated data.

\begin{lemma}[Kernels are Tanner codes of covers]\label{lem:tanner}
\notformalized
Let $N$ be a normal subgroup of $F$ of finite index, and let $X_N\to X_0$ be the regular cover with group $F/N$ and
voltages $\gamma_eN$, with the datum lifted from~$X_0$. Then
\[
  \dim_{\F_3}\ker\varphi_{F/N}(A)=2\dim_{\F_9}T_{\mathrm{tw}}(X_N).
\]
\end{lemma}
\begin{informalproof}
Restriction of scalars commutes with passing to $F/N$, and it doubles dimensions of kernels, so it suffices to show
$\dim_{\F_9}\ker\varphi_{F/N}(A_9)=\dim_{\F_9}T_{\mathrm{tw}}(X_N)$. The vertices of $X_N$ are the pairs $(v,g)$ with
$g\in F/N$, and an edge $e$ of $X_0$ from $v_0$ to $v_1$ lifts to the edges from $(v_0,g)$ to $(v_1,g\gamma_e)$. A row
vector $x\in\F_9[F/N]^{55}$ assigns to each even vertex $(v_0,g)$ of $X_N$ the quartic form $f_{v_0,g}$ whose
coefficients are the coefficients $x_{(v_0,c)}(g)$ of $g$ in $x_{(v_0,c)}$, $c=0,\dots,4$. Let $w$ be the function on the
edges of $X_N$ whose value on an edge is the value of the form at its even end at the label of the edge. For $h\in F/N$,
the coefficient of $h$ in $(x\bar A_9)_{(v_1,r)}$ is
\[
  \sum_{e}m_1(e)_r\sum_{c}m_0(e)_c\,x_{(v_0,c)}(h\gamma_e^{-1})=\sum_{e}m_1(e)_r\,w(\tilde e),
\]
where $e$ runs over the edges of $X_0$ at $v_1$, $v_0$ is the even end of~$e$, and $\tilde e$ is the lift of $e$ with odd
end $(v_1,h)$, whose even end is $(v_0,h\gamma_e^{-1})$. So it is the $r$-th parity check at the odd vertex $(v_1,h)$.
The map $x\mapsto w$ is injective because evaluation of quartics at ten points is injective (\cref{lem:selfdual}), and its
image is the set of edge functions whose restriction to every even vertex lies in the local code. The vectors
$(m_1(e)_r)_e$, $r=0,\dots,4$, span the local code at the odd vertex, which is self-dual (\cref{lem:selfdual}); so the
parity checks at $(v_1,h)$ vanish exactly when the restriction of $w$ to its edges lies in the local code there. Hence
$\ker\varphi_{F/N}(A_9)\cong T_{\mathrm{tw}}(X_N)$.
\end{informalproof}
